% A paragraph ending in \\ (its last line empty) before a theorem, a label and a plain paragraph: the empty line is the paragraph's own, and the space after it is TeX's.
\documentclass[12pt]{article}
\usepackage{amsthm}
\parindent=25pt
\parskip=18pt
\newtheorem{lemma}{Lemma}
% one page with nothing but the text: checked against plain lualatex too (cases.ts)
\pagestyle{empty}
\begin{document}
\baselineskip20pt
Case A, label: a paragraph long enough to fill a line of the text width and then some more words, it ends in a newline.\\
\begin{lemma}\label{a}
Let the standing assumption hold true. Then the relation is symmetric.
\end{lemma}
Case B, no label: a paragraph long enough to fill a line of the text width and then some more words, it ends in a newline.\\
\begin{lemma}
Let the standing assumption hold true. Then the relation is symmetric.
\end{lemma}
Case C, a plain paragraph after: a paragraph long enough to fill a line of the text width and then some, it ends in a newline.\\

A plain paragraph follows, long enough to fill one line of the text width and then some more words here.

Case D, no newline: a paragraph long enough to fill a line of the text width and then some more words, it ends plainly.
\begin{lemma}\label{d}
Let the standing assumption hold true. Then the relation is symmetric.
\end{lemma}
The end.
\end{document}
